% Generated draft: requires mathematical and source review.
\documentclass[11pt]{article}
\usepackage[margin=25mm]{geometry}
\usepackage{fontspec}
\setmainfont{Latin Modern Roman}
\usepackage{amsmath,amssymb,mathtools,unicode-math}
\setmathfont{Latin Modern Math}
\usepackage{enumitem,needspace,xurl,hyperref}
\setlength{\parindent}{0pt}
\setlength{\parskip}{5pt}
\setlength{\emergencystretch}{3em}
\newcommand{\N}{\mathbb{N}}
\newcommand{\Z}{\mathbb{Z}}
\newcommand{\Q}{\mathbb{Q}}
\newcommand{\R}{\mathbb{R}}
\newcommand{\C}{\mathbb{C}}
\newcommand{\F}{\mathbb{F}}
\newcommand{\A}{\mathbb{A}}
\newcommand{\PP}{\mathbb{P}}
\newcommand{\E}{\mathbb{E}}
\newcommand{\eps}{\varepsilon}
\newcommand{\dd}{\,\mathrm{d}}
\begin{document}
\section*{UnsolvedMath ID: 20003187}
\textbf{Source problem code:} catalogue entry
\subsection{Definitions and mathematical statement}
Oscar Garc\'ia-Prada's source question proposes a degree-four Toledo invariant for representations into a real group with quaternion-K\"ahler symmetric space. A precise representation invariant needs more than an arbitrary four-form. For the supplied interpretation, let \(M\) be closed, oriented, and smooth of real dimension four, let \(G\) have a contractible noncompact-type quaternion-K\"ahler symmetric space \(X=G/K\), and choose a closed \(G\)-invariant four-form \(\Omega\) with fixed normalization. Given \(\rho:\pi_1(M)\to G\), choose a \(\rho\)-equivariant map \(f:\widetilde M\to X\). Its pullback descends to \(M\), and
\[
 \tau_\Omega(\rho)=\int_M f^*\Omega.
\]
Here \(f\) is equivariant on the universal cover, not an ordinary map \(M\to X\); contractibility gives the intended independence of the section choice. Compact-type targets or other surface/boundary conventions need a separate formulation.

Let \(Q_\rho\) be the oriented rank-three real bundle of quaternionic structures obtained from the canonical \(\mathrm{SO}(3)\)-bundle over \(X\) and the associated flat \(X\)-bundle/section. The aid's integral normalization is
\[
 \tau_{\mathrm{int}}(\rho)=\langle p_1(Q_\rho),[M]\rangle\in\Z,
\]
with the corresponding Chern--Weil normalization of \(\Omega\). Its Dold--Whitney mod-four constraint is
\[
 p_1(Q_\rho)\bmod4=\mathcal P(w_2(Q_\rho))\in H^4(M;\Z/4\Z),
\]
where \(\mathcal P\) is the Pontryagin square. Evaluation gives the same congruence for the integral number. The canonical bundle choice and form scale must be fixed before using this identity as a numerical normalization.

For quaternionic hyperbolic targets, the aid invokes Pieters's norm theorem to give, in its metric normalization,
\[
 |\tau_{\mathrm{met}}(\rho)|\leq v_4\,\|M\|_1,
\]
where \(\|M\|_1\) is simplicial volume and \(v_4\) is the norm constant in that theorem. For a closed real-hyperbolic four-manifold, using compatible curvature and simplex-volume normalization, it reports
\[
 |\tau_{\mathrm{met}}(\rho)|\leq\operatorname{Vol}(M).
\]
The integral Pontryagin form and metric quaternionic form may differ by a scalar; no equality between \(\tau_{\mathrm{int}}\) and \(\tau_{\mathrm{met}}\) is assumed without that conversion. The target metric, exact \(v_4\) convention, and conversion factor are not supplied.

The explicit integral family uses \(M=\Sigma_g\times\Sigma_h\), \(g,h\geq2\), with product orientation, and factor representations into \(\mathrm{SU}(1,1)\) followed by the block inclusion
\[
 \mathrm{SU}(1,1)\times\mathrm{SU}(1,1)\longrightarrow\mathrm{SU}(2,2).
\]
Let \(L_g,L_h\) be their negative-line bundles, pulled back to the product, with Chern numbers \(d_g,d_h\) on the respective factors. The aid identifies \(Q_\rho=\mathfrak{su}(W)\), \(W=L_g\oplus L_h\), and uses
\[
 p_1(\mathfrak{su}(W))=c_1(W)^2-4c_2(W)
\]
to obtain
\[
 \tau_{\mathrm{int}}(\rho)=-2d_gd_h.
\]
For consistently oriented uniformizing factors it gives
\[
 \tau_{\mathrm{int}}(\rho)=-2(g-1)(h-1)=-\chi(M)/2,
\]
up to orientation conventions, and states compatibility with the Pontryagin-square congruence. This formula is in the integral characteristic-class normalization and is not automatically the metric-normalized value in the preceding inequality.

The original requests remain broader: K\"ahler surfaces are highlighted by Garc\'ia-Prada; Dieter Kotschick suggests aspherical \(M\); Domingo Toledo suggests complex-hyperbolic surfaces, constant-curvature four-manifolds, and products of surfaces. The characteristic constraint, quaternionic-hyperbolic bound, and block-product calculation are supplied partial results, not a classification or a complete theory for all these cases.
\subsection{Short English statement}
The supplied refinement defines the invariant through equivariant pullback, identifies an integral version with a Pontryagin number, and records a mod-four constraint. A metric-normalized volume bound and an explicit product-of-surfaces formula are separate partial results whose normalization scales must be matched.
\subsection{Sources}
\begin{itemize}
\item Local UnsolvedMath registry, permanent ID \(20003187\), exact source code \texttt{catalogue entry}; canonical \texttt{aim-topology-notes.json}, \texttt{notes[274]}, source tag \texttt{question}.
\item AIM workshop ``Representations of surface groups,'' Question 4.2 by Oscar Garc\'ia-Prada, Comment 4.3 by Dieter Kotschick, and Comment 4.4 by Domingo Toledo: \url{https://aimath.org/WWN/surfacegroups/surfacegroups.pdf}. The supplied verification calls the canonical question exact. The PDF was not fetched.
\item The research summary names Pieters's norm theorem and the Dold--Whitney constraint, but provides no titles, bibliographic identifiers, full authoritative citations, or normalization conversion. The original literature field is empty.
\item Research attempt 1: result \texttt{partial\_result}, problem status at run \texttt{partially\_solved}, rigor label \texttt{proved}, novelty confidence \texttt{low}; classification \texttt{PARTIAL-PROGRESS}, registry status \texttt{partially\_solved}. Created and updated 14 August 2026. These are supplied assessments requiring review.
\item Attribution: Ulam AI / UnsolvedMath Contributors, CC BY 4.0. The machine-generated research aid is not peer reviewed and requires independent verification before citation.
\end{itemize}
\end{document}
